Сначала прочитайте число $N$.

Чтобы сделать запрос, выведите <<? x>> $(1 \le x \le N)$ и прочитайте ответ:

\begin{itemize}
\item <<YES>>, если число $x$ любимое;
\item <<NO>>, если число $x$ нелюбимое.
\end{itemize}

Про каждое число можно спросить не более одного раза, порядок запросов любой. Чтобы дать окончательный ответ, выведите <<! F>>, где $F$ --- количество любимых чисел, и завершите программу.

Ваша программа получит вердикт <<Неправильный ответ>> в случае:

\begin{itemize}
\item нарушения формата вывода;
\item запроса числа вне промежутка $[1; N]$;
\item повторного запроса про одно и то же число;
\item неверного окончательного ответа.
\end{itemize}

\textbf{После каждого выведенного запроса обязательно сбрасывайте буфер вывода}, иначе ваша программа получит вердикт <<Idleness Limit Exceeded>>:

\begin{itemize}
\item C++: \texttt{std::endl} при выводе через \texttt{iostream} (\texttt{cout << ... << endl}), \texttt{fflush(stdout)} после вывода через \texttt{stdio} (\texttt{printf});
\item Java: \texttt{System.out.println(...)}, а при выводе через другой поток (\texttt{PrintWriter}, \texttt{BufferedWriter}) --- его \texttt{flush()} после каждой строки;
\item Python: \texttt{print(..., flush=True)}, а ответы подсистемы читайте через \texttt{input()}.
\end{itemize}
